\NewPage{6.1 核心洞察与趋势展望}{把四类证据带回共同任务}{支持持续参与，维护可复用资源，检验公共结果}{s61}

前五章从参与分布、主体影响、技术分工和模型生态观察开放协作。它们共同提示：\key{生态规模、网络位置与开放条件需要分别判断。}行动应回到具体参与者、资源和责任，才能识别投入应当解决的问题。

\plate{figures/new/01_evidence.pdf}{图6-1　四个研究视角如何进入行动问题}{来源：第二章2.1、第三章3.4、第四章4.1及第五章5.2—5.3。图示为证据与问题的对应，不表示因果、权重或成熟度。}

\topic{参与机会：把入口延伸到持续协作}
第二章中，累计参与基础与当期活跃的增长并不同步。年度汇总无法识别同一批人的留存与角色成长。因此，人才支持还需观察首次参与之后的反馈、持续贡献和维护接续。

\topic{开放基础：把影响位置落实到维护对象}
第三章的项目画像与第四章的功能分工提供调查入口；第五章的声明关系则提示中间版本也可能连接下游。实际支持对象仍需由具体系统的依赖、版本和维护需求确定。

\topic{公共结果：从资源可见走向场景可验}
平台活动、下载和关系声明各有解释范围。开放资源怎样改善真实任务、由谁承担持续成本，需要实施与结果证据。后续建议据此提出支持安排，不预设干预已经有效。

\vfill
\noindent\includegraphics[width=\linewidth]{figures/new/chapter6_dimensions.pdf}
\src{十维标识用于联系发展目标，不构成新分组或综合评分。全章默认主统计期：2025-07-01至2026-06-30；不同窗口随文标明。}

\NewPage{6.1 核心洞察与趋势展望}{发布之后，还需什么条件}{材料可查、权利明确与实际可用，需要各自的证据}{s61-open}

第五章将供给增长、派生复用与开放条件分别观察。对于下一步投入，更值得追问的是：\key{已有资源能否被继续研发，哪些交付缺口可以共同补齐？}完整性与开放性提供了两个不同的检查入口。

\plate{figures/new/02_openness.pdf}{图6-2　高开放组内部，材料完整性仍需单独判断}{来源：第五章5.3的\href{https://github.com/X-lab2017/2026-Global-Open-Source-Development-Report/blob/c51e3441361407b68f10b339d97eff288319b503/data/processed/chapter5/fig-5-3-1_openness_quadrants.csv}{冻结评估汇总}。870样本中，高开放600个，其中高完整117个、低完整483个；图中比例分母为600。精确筛选年窗与逐项身份名单尚未取得，不代表平台总体或主期年度结果。}

\topic{从比例判断，转向具体交付}
这组样本支持把材料与权利分开检查，不能替代每个版本的判断。发布团队需要说明权重、代码、数据来源、训练与评测材料的可得性；采用团队则需结合具体任务核对使用条件、运行环境和维护入口。材料入口存在，也不等于已经完成独立复现。

\topic{未来变化，应留下可检验的问题}
如果供给继续增长，适配与维护是否得到相应支持？如果更多中间版本连接上下游，改进能否回到上游、责任能否接续？如果参与入口继续扩大，新参与者能否进入评审与决策？这些是前文证据引出的后续观察问题，尚不能写成确定的未来轨迹。

\vfill
\ReportNote{第五章的2026-09-13可见仓库年度回溯、模型关系快照与本页开放样本各有边界。本章不将它们拼成同一年度增幅，也不据此排序国家整体能力。}

\NewPage{6.2 政府与国际组织}{让公共支持能够持续}{把投入连接到维护、参与条件与实际结果}{s62}

公共支持可以依据前文证据识别调查对象，再与维护者、采用方和本地伙伴确认需求。优先顺序与投入规模应取决于真实任务和资源约束。

\Action{01}{将公共支持延伸到长期维护}{第三章项目画像与第四章4.1的功能分工提示，应用背后涉及多层能力。\key{资助方联合公共采用方、基金会和维护团队}，先核实共同使用的组件及支持缺口，再约定版本、文档、测试、安全响应和人员接续。预算同时覆盖日常维护与协调，明确支持期限和退出交接。}{约定任务的完成与延期、问题处理及维护角色接续；同时记录投入。}

\Action{02}{将人才与国际合作连接到真实任务}{第二章区分参与基础与当期活跃，却尚未解释同批次参与者的成长。\key{政府与国际组织联合高校、企业和社区}，围绕文档、测试、数据整理、部署适配或评测提供导师与协作资源。让全球南方及其他资源较少地区的伙伴共同定义任务、交付和评价。}{固定参与批次的持续贡献与角色变化，另记导师投入、退出与失访。}

\Action{03}{分开验收交付、采用与公共结果}{第五章显示材料与开放条件不能互相替代。\key{公共采购或项目支持方}应进一步分开保存开放交付、真实运行与结果评价三类证据。与本地使用者先确定服务问题、现有做法和成本，再检验实施后的变化；不能以交付数量或下载量替代实际成效。}{材料能否核查、系统能否运行、目标任务是否改善，以及新增负担。}

\vfill
\ReportNote{证据入口：第二章2.1；第三章3.4；第四章4.1；第五章5.3；第一章公共能力条件关系。三条建议均为有条件的行动安排，不设统一国家处方、资助金额或成效阈值。}

\NewPage{6.2 政府与国际组织}{一项公共投入，三层验收}{先明确问题与起点，再判断交付、采用和结果}{s62-evaluation}

公共项目最容易取得的是交付记录，最终需要回答的却是使用者发生了什么变化。把三层证据分开，才能识别问题停在哪里，以及下一轮支持应补什么。

\plate{figures/new/03_evaluation.pdf}{图6-3　开放交付、实际采用与公共结果分别验收}{依据第一章公共能力框架、第三章3.4与第五章5.3归纳的实施建议。箭头表示检查顺序；每层允许“未验证”或“不适用”，上一层完成不表示下一层已经实现。}

\topic{从一个真实场景开始}
由使用者、实施方和支持方共同确定服务对象与任务，记录现有做法、一次性适配和持续运维成本。结果指标围绕场景选取，例如任务完成、错误与返工、服务可达性；没有可靠起点和比较设计时，宜描述观察变化，不直接作因果判断。

\topic{把继续、调整与退出都纳入安排}
材料可查而运行失败，应定位环境、兼容、人员或流程问题；运行稳定而结果未改善，应重新检查任务和支持条件。同步记录失败、额外负担与不利影响，为调整、缩小实施或有序退出保留依据。

\vfill
\ReportNote{公共价值需要场景验证。SDGs及数字公共品参照可以帮助确定评价问题，不构成对具体项目的认证或成效证明。}

\NewPage{6.3 企业与研究机构}{把使用、发布与复核连起来}{承担实际依赖的维护成本，让开放交付可以检查}{s63}

企业和研究机构既使用开放资源，也发布成果并影响评价方式。行动重点是将资源交付落实到具体版本，将持续使用落实到投入和责任。

\Action{04}{沿实际依赖承担维护责任}{第三章项目画像与第四章功能全景提供线索，但分类相邻不等于实际依赖。\key{采用方的工程、产品、采购与开源治理团队}应核实运行环境中的版本与责任，约定兼容测试、升级、安全响应和替代路径。对关键上游，以工程贡献、维护合同或资助承担合理工作量。}{投入是否对应约定维护任务；结合问题难度查看升级、响应与故障处理。}

\Action{05}{按具体版本提供可检查的开放交付}{第五章的开放样本表明，许可与材料需要分别改进。\key{发布团队}应逐项列明代码、权重、数据说明、训练与评测材料、使用条件和维护入口。无法公开的材料说明缺失及原因；发布后保留变更记录，避免用标签或规模概要代替实际材料。}{分别记录“入口存在”“材料已核对”“独立运行通过”，保留未知项。}

\Action{06}{把研究支持延伸到独立复核}{第二至第五章的身份、窗口与样本差异限制了比较。\key{高校、研究机构与平台}应支持数据维护、文档、测试和独立复核，优先修复会改变结论的证据缺口。保留方法版本、筛选与缺失处理；受限制的材料提供适当的授权复核路径。}{复核对象与环境、成功范围、失败原因及维护投入；不只看下载和榜位。}

\vfill
\ReportNote{证据入口：第三章3.4；第四章4.1及全景中的对象边界；第五章5.3。支持对象须由实际依赖确认，不把4.6的PyTorch生态观察样本视为基金会正式托管名单。}

\NewPage{6.3 企业与研究机构}{从版本材料进入真实任务}{两项已有观察，转成采用与复核的具体动作}{s63-cases}

第五章的固定版本案例提供了材料核查入口。本章进一步提出任务：\key{将材料检查推进到环境、适配、维护与使用结果。}这一步需要实际实施记录，不能仅凭模型卡完成。

\plate{figures/new/04_cases.pdf}{图6-4　已有版本材料支持什么行动}{来源：第五章5.3及\href{https://huggingface.co/allenai/OLMo-2-1124-7B/blob/7df9a82518afdecae4e8c026b27adccc8c1f0032/README.md}{OLMo固定模型卡}、\href{https://huggingface.co/meta-llama/Llama-3.1-8B/blob/d04e592bb4f6aa9cfee91e2e20afa771667e1d4b/README.md}{Llama固定模型卡}。沿用2026-10-17材料记录及2026-11-07补充核读，均为期后材料观察；外链代码与数据未全部固定于同一提交。两例不并入870样本，不重新授级。}

\topic{试点需要同时留下四类记录}
\key{成本：}分别记录一次性适配、持续运行、维护和人员协调。\key{本地化：}说明语言、数据、任务与工作流程的适配范围。\key{维护：}落实责任人、支持入口、升级与接续安排。\key{结果：}观察实际任务表现、使用者收益与新增负担。

\topic{用小范围验证决定是否扩大}
采用方先选择一个真实任务和固定版本，完成环境记录、材料核对与约定测试；研究方保留成功及失败边界，维护团队确认新增工作量。再依据这些记录判断能否扩大使用，并明确还缺哪些支持。

\vfill
\ReportNote{两例尚未覆盖完整训练复现及成本、本地化、维护与结果的全部证据。本页提炼核查动作，不将材料可得性升级为实施成功或普惠成效。}

\NewPage{6.4 社区与全球伙伴}{让参与进入维护与决策}{开放协作入口，也为承担责任的人提供支持}{s64}

社区连接贡献者、上游维护者与本地使用者。持续合作既需要透明入口，也需要可承受的工作量、资源安排和反馈机制。

\Action{07}{将首次贡献连接到维护角色}{第二章的参与统计尚未说明谁持续留下、如何成长。\key{社区与支持伙伴}可从边界明确的任务开始，提供文档、导师反馈和异步协作，再公开评审、发布与维护者接续规则。代码、测试、翻译和用户支持均可成为入口，培养成本不应全部由志愿者承担。}{固定批次中实际进入评审或维护角色的记录，另记贡献类型、投入与退出障碍。}

\Action{08}{围绕共同需求组织跨地区维护}{第五章的声明关系帮助寻找上下游伙伴，但不能代替合作事实。\key{上游、下游与本地团队}应共同选择兼容测试、语言适配、文档或评测任务，约定投入、成果使用条件、上游回流与后续支持。资源较少地区的伙伴应参与需求定义和决策。}{共同交付、上游处理状态与跨版本维护，同时观察本地持续使用和改进能力。}

\Action{09}{让使用者反馈进入治理}{第一章把公共价值视为需要验证的目标。\key{社区、实施方与使用者代表}可建立需求、变更和反馈入口，公开重要决策及理由，说明哪些意见改变了版本或安排。支持机构应承担必要协调和复核，避免将问责与记录负担继续压给维护者。}{谁能参与、哪些需求进入决策、谁仍被遗漏，以及反馈后的实际变化。}

\vfill
\ReportNote{证据入口：第二章2.1；第五章5.2；第一章公共价值框架。观察应优先使用任务所需的汇总记录，不把暂停参与当作个人失败，也不以单一提交量评价长期合作。}

\NewPage{6.5 结语与后续观察}{让行动留下可检验的经验}{下一轮优先连接参与、资源与结果的连续记录}{s65}

本报告用四个视角识别开放生态的不同侧面。下一步需要把这些观察连接到真实使用：谁持续参与，哪些资源得到维护，哪些任务发生改善。\key{尚未连接起来的证据，正是后续观察的重点。}

\plate{figures/new/05_observation.pdf}{图6-5　下一轮优先建立的三条连续记录}{依据本章三项任务与九条建议归纳。箭头表示拟建立的记录联系，不代表已经实现的因果关系；三条记录分别保留对象、窗口与责任，不合成统一总分。}

\topic{先确定观察对象，再比较变化}
参与观察固定批次，资源观察固定版本与依赖，结果观察固定任务和服务对象。范围变化、退出、失访和材料缺失分别记录，不能把不可观察当作零。数量与比例需要和工作难度、投入及具体经历一起解释。

\topic{让各方承担相应责任}
公共支持改善基础条件，采用方承担持续使用的成本，发布者提供可检查的材料，研究机构开展独立复核，社区与本地伙伴共同决定需求和接续。合作能否延续，需要明确资源来源，不能寄托于另一方无限增加劳动。

\vfill
\ReportNote{本章提出后续观察安排，尚未测得九条建议的干预效果。全书最终关注的是：更多人能否持续参与，开放成果能否被维护与再创造，以及这些工作是否形成可验证、可延续的公共价值。}
